\begin{proposition}[ARI \ref{ari:C.6.R2}; Proved: central derivatives with explicit tails]\label{ari:C.6.R2}
The following are the values of $L'(E_i,1)/\Omega_i$; zeros in rows
one and eight are exact. More digits, the periods, derivatives and
tail bounds are in \path{results_Crank.csv}.
\begin{center}\small
\begin{tabular}{r|l}
$i$&$L'(E_i,1)/\Omega_i$\\\hline
1&$0$\\
2&$5.134398425737889664345619204085016979705061236171249722654226\ldots$\\
3&$2.134120849044754020051650801945664635106456422166881627967598\ldots$\\
4&$31.81013979419927776012209427941545866492237285892013668350005\ldots$\\
5&$1.869462708725885120030080805051387134542481374813232701178304\ldots$\\
6&$3.428873379050567976208155749688133688667614469824389504923957\ldots$\\
7&$0.362583348071598548315763082767526599722077924892872839203426\ldots$\\
8&$0$\\
9&$9.629273918234131678053742152263035417231796757756973886621556\ldots$\\
10&$16.04238857830583792513538634637663343755958963581574670348088\ldots$\\
11&$1247.321806154531180413038860827393095408344267114181914400926\ldots$
\end{tabular}
\end{center}
The two second derivatives, divided by the same real periods, are
\begin{align*}
 L''(E_1,1)/\Omega_1
 &=28.71408722758099118128474692322079946598060512912402172882525\ldots,\\
 L''(E_8,1)/\Omega_8
 &=78.03222041637032223732443782051140882655135939512581188900823\ldots.
\end{align*}
In particular the leading Taylor coefficients for these two curves
are half the displayed second derivatives.
\end{proposition}
\Uses{the hypotheses, notation and definitions stated here; \cref{ari:eq:Cr-tails}; \citet[Section 3]{CrDarmon2009}; \citet[Theorem 10.1]{CrRubin1999}}

\begin{proof}[Derivative formulas and error control]
\begin{enumerate}[label=\arabic*.,leftmargin=*]
\item Let $N$ be the conductor, $c=2\pi/\sqrt N$, and
$L(E,s)=\sum_{n\geq1}a_nn^{-s}$. Put
\[
 E_1(x)=\int_x^\infty\frac{e^{-u}}u\,du,
 \qquad G(x)=\int_x^\infty e^{-u}\frac{\log(u/x)}u\,du.
\]
Splitting the Mellin integral at $1/\sqrt N$ and using the functional
equation gives, for root number $-1$,
\begin{equation}\label{ari:eq:Cr-first}
 L'(E,1)=2\sum_{n\geq1}\frac{a_n}{n}E_1(cn).
\end{equation}
For root number $+1$ and $L(E,1)=0$, it gives
\begin{equation}\label{ari:eq:Cr-second}
 L''(E,1)=4\sum_{n\geq1}\frac{a_n}{n}G(cn).
\end{equation}
One can verify the factors directly by writing the completed Mellin
integral as an integral over $[1,\infty)$ of
$(u^{s-1}+wu^{1-s})f(iu/\sqrt N)$: its first derivative has
the factor $2\log u$ when $w=-1$, and its second derivative
the factor $2\log^2u$ when $w=1$. Integration by parts gives
the two formulas above. The vanishing of the lower derivatives
removes all derivatives of the gamma and conductor factors.

\item Hasse's bound and multiplicativity imply
$|a_n|\leq d(n)\sqrt n\leq2n$. Also
$E_1(x)\leq e^{-x}/x$ and $G(x)\leq e^{-x}/x^2$.
The latter follows on writing $u=x+v$ and using
$\log(1+v/x)\leq v/x$. Consequently the absolute tails after
$n=M$ are at most
\begin{equation}\label{ari:eq:Cr-tails}
 \frac{4e^{-c(M+1)}}{c(M+1)(1-e^{-c})},\qquad
 \frac{8e^{-c(M+1)}}{c^2(M+1)^2(1-e^{-c})},
\end{equation}
respectively. For the first ten curves we used
$M=\lceil180/c\rceil$, exact integer coefficients, and 1024-bit
MPFR arithmetic. For the rank-two kernel we used the convergent series
\[
 G(x)=\frac12(\log x+\gamma)^2+\frac{\pi^2}{12}
       +\sum_{k\geq1}\frac{(-x)^k}{k^2k!}.
\]
After $k>x$ its terms decrease in absolute value. The alternating
remainder was cut below $10^{-140}$. At most 740 terms were needed;
the largest intermediate term was less than $1.4\cdot10^{72}$.
The 1024-bit arithmetic therefore leaves far more than 100 digits
after this cancellation. The largest bound in
\eqref{ari:eq:Cr-tails} was $3.46\cdot10^{-77}$. The source and complete
outputs are \path{code/Crank_coefficients.gp},
\path{code/Crank_derivatives.cpp}, and their corresponding outputs.

\item For $i=11$, a full vector of $819$ million coefficients would be
unnecessarily large. We instead evaluated the CM theta series in
bounded-memory segments. The curve is the quadratic twist, by
$d=185801=7\cdot11\cdot19\cdot127$, of the conductor-$163^2$
curve
\[
 y^2+y=x^3-2174420x+1234136692.
\]
Write $D=163$, $\chi_D(a)=(a/D)$, and $\chi_d(n)=(d/n)$.
The exact coefficient $a_n$ is the sum of
\[
 \chi_d(n)\chi_D(a)a \quad(n=a^2,\ a>0)
\]
and
\[
 \chi_d(n)\chi_D((D+1)u/2)u
 \quad\left(n=\frac{u^2+Dv^2}{4},\ u,v>0,
                           \ u\equiv v\pmod2\right).
\]
This is the Hecke-character theta series for
$\mathcal O_{\Q(\sqrt{-D})}$, grouping conjugate algebraic integers
and then applying the quadratic twist; see
\citet[Section 3]{CrDarmon2009}. The first 10000 coefficients agree
exactly with PARI for both the untwisted and twisted curve. The same
code at $D=67,d=217$ reproduces the independent PARI derivative.

\item To accelerate the sum, in a block $n=n_0+r$, $|r|\leq R$, form the
exact integer moments $M_k=\sum a_nr^k$. For $x=cn_0$ write
\[
 e_k=e^{-x}(-x)^k/k!,\quad g_0=e^{-x},\quad
 g_k=e_k-g_{k-1},\quad
 \phi_0=E_1(x),\quad\phi_k=-\phi_{k-1}-g_{k-1}/k.
\]
Then the block contribution to $L'/2$ is
$\sum_k\phi_k M_k/n_0^{k+1}$. This follows by Taylor expansion
of $E_1(x(1+s))/(1+s)$. On $|s|=1/2$ its modulus is at most
$2E_1(x/2)$, which is less than $14$ throughout the block range.
We take enough terms that $(2R/n_0)^{k+1}<10^{-85}$.
The Cauchy remainder, summed with the actual
$\sum|a_n|<2\cdot10^{12}$ and $n_0>7700$, is less than
$10^{-75}$. The values of $E_1$ at consecutive block centres
are continued by the analogous Taylor expansion and refreshed
from a correctly rounded MPFR evaluation every 1024 blocks.
The continuation bound is $10^{-110}$ per normalized expansion.
All moments are integers, not floating-point sums.

\item The 384-bit run used $R=128$ and cutoff $170$; a separate 512-bit
run used $R=96$ and cutoff $190$. Their values agree beyond 75
decimal places. The latter used $915818441$ coefficients in
segments, with an omitted tail less than $3.10\cdot10^{-78}$,
and gave
\[
 L'(E_{11},1)=
 2.296582198163380608965915774436446055185027904705664673866754\ldots.
\]
The elementary operation-rounding bounds at these precisions are
smaller than $10^{-90}$: in the recurrences above
$|e_k|\leq1$, $|g_k|\leq1$, and $|\phi_k|<20$;
there are fewer than $5\cdot10^6$ blocks and the recurrences are
refreshed before 1024 consecutive continuations. Bounding each
normalized moment by $\sum|a_n|(R/n_0)^k$ gives this estimate
without assuming cancellation in the total sum. The arithmetic
and truncation margins are amply sufficient for every displayed
50-digit value. Both executable builds and both transcripts are
retained as \path{code/Crank_cm_stream*}.

\item The real periods are enclosed from the exact minimal invariants,
independently of the saved decimal period values.  For rows $1,2,3,4$
put respectively $T=\sqrt2,\sqrt3,2,\sqrt7$ and $\tau=iT$;
for rows $5,\ldots,11$ put
$T=\sqrt{7}/2,\sqrt{11}/2,\sqrt{19}/2,\sqrt{27}/2,
\sqrt{43}/2,\sqrt{67}/2,\sqrt{163}/2$ and $\tau=1/2+iT$.
Let $r=e^{-2\pi T}$ and $q=e^{2\pi i\tau}$.  The identity
$E_4=1+240\sum_{n\geq1}\sigma_3(n)q^n$ and
$\sigma_3(n)<2n^3$ give, after $M=100$ terms, the tail bound
\[
480r^{A}\left\{
\frac{A^3}{1-r}+\frac{3A^2r}{(1-r)^2}
+\frac{3Ar(1+r)}{(1-r)^3}
+\frac{r(1+4r+r^2)}{(1-r)^4}\right\},\qquad A=101.
\]
This follows by expanding $(A+j)^3$ and summing each geometric
moment.  All eleven parameters satisfy
$1-1008r/(1-32r)>0$; the bound
$\sigma_5(n)<2n^5$ and $n^5\leq32^{n-1}$ therefore proves
$E_6(\tau)>0$.  The exact Weierstrass invariant relations yield
$\kappa=2\pi(E_4/c_4)^{1/4}$.  For $c_6>0$ the real lattice
generator has length $\kappa$.  For $c_6<0$ the lattice is rotated
by $i$: its real generator has length $T\kappa$ in the rectangular
case and $2T\kappa$ in the other case.  Indeed the real vectors in
$i(\Z+iT\Z)$ are $T\Z$, while those in
$i(\Z+(1/2+iT)\Z)$ are $2T\Z$.  Multiplication by the number
$n_{\R}=2$ for $\Delta>0$ and $n_{\R}=1$ for $\Delta<0$
gives the full real period used in the statement.

The current program \path{code/I/certify_all_period_ratios.py}
performs these operations with 768-bit Arb intervals and verifies
every displayed period and derivative ratio.  For the first ten
derivatives it enlarges the explicit tail by one percent and adds
$10^{-100}$ for arithmetic and kernel-series errors.  This last
bound is conservative: fewer than $5\cdot10^5$ coefficients and
740 series terms occur, the intermediate terms are below
$1.4\cdot10^{72}$, and the 1024-bit accumulated rounding and
$10^{-140}$ kernel remainders are below this allowance.
For the last derivative it adds $1.1\cdot10^{-75}$, covering the
proved block Taylor and rounding bounds above.  Division by the
strictly positive period intervals preserves all fifty displayed
digits.  The resulting intervals and exact invariants are in
\path{code/I/records/all_period_ratios_arb.json}.

\item The exact root numbers are $+1$ for $i=1,8$ and $-1$ otherwise.
All points are nontorsion by Statement C.5. For a CM curve,
the Coates--Wiles theorem, in the form
\citet[Theorem 10.1]{CrRubin1999}, says that a nonzero central
Hecke $L$-value forces the Mordell--Weil group over the CM field
to be finite. Proposition~5.3 of the same source supplies, for a nonmaximal CM order, an isogenous
curve with maximal order: the quotient is by the Galois-stable subgroup attached to the conductor ideal. The isogeny induces an isomorphism on Mordell--Weil groups tensored with $\Q$ and on their rational Tate modules, preserving ranks and local $L$-factors.
Since $L(E/K,s)=L(E/\Q,s)^2$ for the CM field $K$, nontorsion
therefore forces $L(E,1)=0$. This proves the needed central
vanishing also in the two sign-$+1$ cases. The functional equation
then forces $L'(E,1)=0$ for those two curves. The positive first
or second derivatives just evaluated prove exactly the asserted
analytic ranks.
\end{enumerate}
\end{proof}